\documentclass{article}

\usepackage[T1]{fontenc}
\usepackage{lmodern}
\usepackage{amsmath}
\usepackage{amssymb}
\usepackage{amsthm}
\usepackage{mathtools}
\usepackage[hidelinks]{hyperref}

\newtheorem{lemma}{Lemma}
\newtheorem*{mainstatement}{Main statement}

\begin{document}

\section{Exact separation of an integer period lattice}\label{sec:title}

\textbf{Draft status: PROPOSITION, pending independent hostile review and Lean replay.}

\begin{mainstatement}
  Let \(u, v\) be integers, not both zero, such that \(281\) divides \(7u + 145v\). Then
  \(\lvert u\rvert + \lvert v\rvert \ge 23\). Equality is attained at
  \((u, v) = (9, -14)\). The assertion concerns only this explicitly defined integer
  lattice.
\end{mainstatement}

\section{Definitions}\label{sec:definitions}

For an integer \(t\), \(\lvert t\rvert = t\) when \(t \ge 0\), and \(\lvert t\rvert = -t\)
when \(t < 0\). For integers \(m, t\), the notation \(m \mid t\) means that \(t = mk\) for
an integer \(k\). The lattice in the statement is the set
\begin{equation}
  \Lambda = \{(u, v) \in \mathbb{Z}^2 : 281 \mid (7u + 145v)\}.
\end{equation}
The length used here is the sum \(\lvert u\rvert + \lvert v\rvert\), not Euclidean length.
A basis means that every lattice vector has a unique expression as an integer linear
combination of the displayed vectors.

\section{An explicit basis}\label{sec:basis}

\begin{lemma}\label{lem:basis}
  The vectors \(p = (9, -14)\) and \(q = (13, 11)\) are a basis of \(\Lambda\). For
  \((u, v) \in \Lambda\), writing \(7u + 145v = 281k\), its coefficients are
  \begin{equation}
    a = 122k - 3u - 63v, \quad b = 2k - v,
  \end{equation}
  and satisfy
  \begin{gather}
    281a = 11u - 13v, \quad 281b = 14u + 9v, \\
    u = 9a + 13b, \quad v = -14a + 11b.
  \end{gather}
\end{lemma}

\begin{proof}
  First, \(7 \cdot 9 + 145 \cdot (-14) = 63 - 2030 = -1967 = -7 \cdot 281\) and
  \(7 \cdot 13 + 145 \cdot 11 = 91 + 1595 = 1686 = 6 \cdot 281\). Thus integer
  combinations of \(p, q\) belong to \(\Lambda\). Conversely, choose the integer \(k\) in
  the definition of divisibility and use the formulas for \(a, b\). Substitution gives
  \begin{align}
    281a &= 122(7u + 145v) - 843u - 17703v \\
         &= (854 - 843)u + (17690 - 17703)v = 11u - 13v,
  \end{align}
  \begin{equation}
    281b = 2(7u + 145v) - 281v = 14u + 9v.
  \end{equation}
  Consequently
  \begin{gather}
    281(9a + 13b) = 9(11u - 13v) + 13(14u + 9v) = 281u, \\
    281(-14a + 11b) = -14(11u - 13v) + 11(14u + 9v) = 281v.
  \end{gather}
  Cancellation of the nonzero integer \(281\) proves the reconstruction formulas. For any
  other representation \(u = 9a' + 13b'\) and \(v = -14a' + 11b'\), expansion gives
  \(11u - 13v = (99 + 182)a' + (143 - 143)b' = 281a'\) and
  \(14u + 9v = (126 - 126)a' + (182 + 99)b' = 281b'\). Comparing these identities with
  those for \(a, b\) and cancelling \(281\) gives \(a' = a\) and \(b' = b\). This proves
  uniqueness.
\end{proof}

\section{Small length forces small basis coefficients}\label{sec:small}

\begin{lemma}\label{lem:small}
  If \((u, v) \in \Lambda\) and \(\lvert u\rvert + \lvert v\rvert \le 22\), its
  coefficients in Lemma~\ref{lem:basis} belong to \(\{-1, 0, 1\}\).
\end{lemma}

\begin{proof}
  The definition of absolute value gives \(-\lvert t\rvert \le t \le \lvert t\rvert\).
  Adding these inequalities for two integers \(s, t\) gives
  \(-\lvert s\rvert - \lvert t\rvert \le s + t \le \lvert s\rvert + \lvert t\rvert\),
  hence \(\lvert s + t\rvert \le \lvert s\rvert + \lvert t\rvert\): when \(s + t \ge 0\)
  use the upper bound, and when \(s + t < 0\) negate the lower bound. These two cases
  exhaust the integers. For a nonnegative integer \(c\),
  \(\lvert ct\rvert = c\lvert t\rvert\), using the cases \(t \ge 0\) and \(t < 0\) in the
  definition. Also \(\lvert -t\rvert = \lvert t\rvert\) by those same two cases.
  Therefore Lemma~\ref{lem:basis} implies
  \(281\lvert a\rvert = \lvert 11u - 13v\rvert \le 11\lvert u\rvert + 13\lvert v\rvert
  \le 13(\lvert u\rvert + \lvert v\rvert) \le 286 < 562 = 2 \cdot 281\),
  \(281\lvert b\rvert = \lvert 14u + 9v\rvert \le 14\lvert u\rvert + 9\lvert v\rvert
  \le 14(\lvert u\rvert + \lvert v\rvert) \le 308 < 562 = 2 \cdot 281\).
  Thus \(\lvert a\rvert < 2\) and \(\lvert b\rvert < 2\). An integer with absolute value
  less than two is at least \(-1\) and at most \(1\), and the only integers in that
  interval are \(-1, 0, 1\). This proves the list of possibilities without a search.
\end{proof}

\section{Proof of the main statement}\label{sec:main-proof}

Suppose its lower bound fails. The length is an integer, so it is at most \(22\).
Lemmas~\ref{lem:basis} and~\ref{lem:small} apply. The following table lists \textbf{all}
possibilities: each of \(a\) and \(b\) belongs to the three-element set proved in
Lemma~\ref{lem:small}, and for each value of \(a\) all three values of \(b\) appear. Its
third column substitutes into \((u, v) = (9a + 13b, -14a + 11b)\); its last column adds
the two absolute values.

\begin{center}
  \begin{tabular}{llll}
    \(a\) & \(b\) & \((u, v)\) & \(\lvert u\rvert + \lvert v\rvert\) \\
    \hline
    \(-1\) & \(-1\) & \((-22, 3)\) & \(25\) \\
    \(-1\) & \(0\) & \((-9, 14)\) & \(23\) \\
    \(-1\) & \(1\) & \((4, 25)\) & \(29\) \\
    \(0\) & \(-1\) & \((-13, -11)\) & \(24\) \\
    \(0\) & \(0\) & \((0, 0)\) & \(0\) \\
    \(0\) & \(1\) & \((13, 11)\) & \(24\) \\
    \(1\) & \(-1\) & \((-4, -25)\) & \(29\) \\
    \(1\) & \(0\) & \((9, -14)\) & \(23\) \\
    \(1\) & \(1\) & \((22, -3)\) & \(25\) \\
  \end{tabular}
\end{center}

Every nonzero row has length greater than \(22\). The sole remaining row is the zero
vector, excluded by hypothesis. This proves the lower bound. Finally \(9 + 14 = 23\) and
\(7 \cdot 9 + 145 \cdot (-14) = -7 \cdot 281\), so the stated nonzero vector attains the
bound. \qed

All displayed arithmetic consists of multiplication, addition, and cancellation of the
stated integers. No generated case list, numerical certificate, or exhaustive lattice
search is an input to this proof.

\end{document}
